\section{Ciberseguridad y formación universitaria}

La ciberseguridad es un área fundamental dentro de la formación universitaria en carreras relacionadas con informática y sistemas, debido a que permite comprender los riesgos asociados al uso de tecnologías digitales. En este capítulo se presentan los conceptos principales de ciberseguridad, sus principios fundamentales, las amenazas y vulnerabilidades más frecuentes, así como su importancia dentro de la educación superior. Estos elementos sirven como base para comprender por qué el aprendizaje de ciberseguridad no debe limitarse a contenidos teóricos, sino que debe complementarse con actividades prácticas orientadas a la toma de decisiones frente a situaciones de riesgo.

\subsection{Concepto de ciberseguridad}

La ciberseguridad comprende el conjunto de principios, procedimientos, prácticas y controles orientados a proteger sistemas, redes, dispositivos, servicios e información frente a accesos no autorizados, alteraciones, interrupciones, ataques digitales o pérdidas de datos. Desde una perspectiva técnica, su propósito se relaciona con la prevención del daño, la protección y la recuperación de sistemas electrónicos y de la información que estos contienen, procurando preservar su confidencialidad, integridad y disponibilidad (National Institute of Standards and Technology [NIST], s. f.).

En el contexto actual, la ciberseguridad no se limita únicamente a proteger infraestructura tecnológica. También involucra a los usuarios, quienes pueden convertirse en un punto de riesgo cuando no reconocen señales de amenaza o cuando toman decisiones inseguras frente al uso de sistemas digitales. Esta situación es relevante porque muchas brechas de seguridad siguen relacionadas con el comportamiento humano, la ingeniería social, el \textit{phishing}, las credenciales robadas, la explotación de vulnerabilidades de software y los ataques de \textit{ransomware} (Verizon, 2026).

Desde una perspectiva académica, la ciberseguridad no debe entenderse únicamente como un contenido técnico aislado, sino como una competencia transversal para los estudiantes de áreas informáticas. Su aprendizaje requiere comprender conceptos básicos, identificar amenazas frecuentes y aplicar medidas de protección en escenarios concretos. Por esta razón, la formación universitaria en esta área debe integrar teoría, análisis de casos y práctica guiada.

\subsection{Principios fundamentales}

La seguridad de la información suele organizarse alrededor de tres principios básicos: confidencialidad, integridad y disponibilidad. La confidencialidad busca que los datos solo puedan ser consultados por usuarios autorizados; la integridad pretende evitar modificaciones no permitidas de la información; y la disponibilidad procura que los sistemas, servicios y datos estén accesibles cuando se los necesita (NIST, s. f.).

Estos principios permiten explicar por qué un incidente digital no afecta únicamente a un sistema aislado, sino que puede interrumpir procesos académicos, administrativos o personales. Por ejemplo, una falla de confidencialidad puede exponer información sensible; una falla de integridad puede alterar datos importantes; y una falla de disponibilidad puede impedir el acceso a servicios necesarios para el funcionamiento de una institución.

En consecuencia, el aprendizaje de ciberseguridad requiere que el estudiante reconozca tanto el funcionamiento de los mecanismos de protección como las posibles consecuencias de una mala decisión. Comprender estos principios permite analizar mejor los escenarios de riesgo y establecer una relación más clara entre la teoría de seguridad y su aplicación práctica.

\subsection{Amenazas y vulnerabilidades}

Las amenazas digitales son acciones, eventos o condiciones que pueden comprometer la seguridad de un entorno informático. Entre las más frecuentes se encuentran el \textit{phishing}, el \textit{malware}, el \textit{ransomware}, la ingeniería social, el robo de credenciales y la explotación de fallas en sistemas o aplicaciones. Estas amenazas pueden aprovechar tanto debilidades técnicas como errores humanos, por lo que su estudio resulta importante dentro de la formación de estudiantes de informática y sistemas (Verizon, 2026).

Las vulnerabilidades, por su parte, son debilidades presentes en sistemas, configuraciones, procedimientos o hábitos de uso que pueden ser aprovechadas por una amenaza. Una contraseña débil, un equipo desactualizado, permisos excesivos o la falta de verificación de enlaces sospechosos pueden convertirse en puntos de entrada para un ataque. Por esta razón, la formación universitaria debe ir más allá de la definición de conceptos y promover la identificación práctica de riesgos (NIST, s. f.).

Los reportes actuales sobre incidentes de seguridad muestran que muchas brechas siguen relacionadas con factores como credenciales comprometidas, errores humanos, ingeniería social y explotación de vulnerabilidades. El \textit{Data Breach Investigations Report} de Verizon analiza de forma anual patrones de incidentes y brechas de seguridad, destacando la persistencia de amenazas asociadas al comportamiento del usuario y al aprovechamiento de debilidades técnicas (Verizon, 2026).

En el contexto de este proyecto, las amenazas y vulnerabilidades no se abordan únicamente como definiciones aisladas, sino como situaciones que pueden ser representadas mediante escenarios de aprendizaje. Esto permite que el estudiante relacione conceptos de ciberseguridad con casos prácticos, como reconocer un intento de \textit{phishing}, identificar una mala práctica de seguridad o seleccionar una respuesta adecuada frente a un posible incidente.

\subsection{Importancia en la educación superior}

En la educación superior, especialmente en carreras de informática y sistemas, la ciberseguridad debe abordarse como una capacidad profesional que articula análisis, toma de decisiones y respuesta ante incidentes. El estudiante no solo necesita saber qué es una amenaza, sino también comprender cómo reaccionar frente a ella, qué medidas aplicar y qué consecuencias pueden derivarse de una acción incorrecta.

Cuando esta formación se desarrolla únicamente de forma teórica, puede producirse una distancia entre el conocimiento conceptual y la aplicación real. Esta situación es especialmente relevante en ciberseguridad, ya que muchas decisiones deben tomarse frente a escenarios dinámicos, donde el usuario debe interpretar señales de riesgo y actuar de forma oportuna.

De ahí que resulten valiosas las estrategias que introducen práctica, simulación y retroalimentación durante el proceso de aprendizaje. En este proyecto, la ciberseguridad se aborda desde una perspectiva formativa, donde el videojuego educativo no constituye el fin principal, sino un medio para representar escenarios de amenaza y fortalecer habilidades prácticas relacionadas con la identificación, análisis y respuesta ante incidentes.
